%% 24장 특이값 분해의 계산
\chapter{특이값 분해의 계산}
\label{ch:svd}

\abstract{특이값 분해는 4장에서 본 대로 행렬의 랭크와 노름과 조건수와 최선 저랭크 근사를 한꺼번에 알려 주는 가장 강력한 분해이다. 이 장은 그것을 계산하는 방법을 다룬다. 특이값은 $A^{\mathsf T}A$의 고유값의 제곱근이지만, 그렇게 계산하면 작은 특이값이 반올림 오차에 묻힌다는 것을 먼저 직접 확인한다. 그래서 $A$ 자체를 직교 변환으로 이중대각 행렬로 줄이고, 그 이중대각 행렬에 23장의 QR 알고리즘을 암시적으로 적용한다. 이중대각 행렬의 특이값은 원소들에 대해 상대적으로 잘 정해진다는 놀라운 성질이 있고, 그 성질을 지키는 알고리즘을 쓰면 아주 작은 특이값까지 높은 상대 정확도로 구할 수 있다. 특이값만 필요할 때 쓰는 dqds 알고리즘과, 특이벡터까지 구할 때 빠른 분할 정복도 다룬다.}

\section{왜 $A^{\mathsf T}A$를 쓰지 않는가}
\label{sec:why-not-ata}

$A=U\Sigma V^{\mathsf T}$이면 $A^{\mathsf T}A=V\Sigma^2V^{\mathsf T}$이므로, $A^{\mathsf T}A$의 고유값 분해에서 특이값과 오른쪽 특이벡터를 얻을 수 있을 것 같다. 그러나 22장의 정규방정식과 같은 이유로 이 길은 정확도를 잃는다. $A^{\mathsf T}A$를 만드는 순간 반올림 오차가 $u\|A\|^2=u\sigma_1^2$ 정도 생기고, 대칭 고유값 문제는 4장의 Weyl 부등식에 따라 고유값을 절대오차 $u\sigma_1^2$ 정도로만 정하므로, 특이값 $\sigma_i$의 상대오차는 대략 $u(\sigma_1/\sigma_i)^2$이다.

특이값이 $1$부터 $10^{-10}$까지 로그 간격으로 놓인 $50\times8$ 행렬에서 두 방법을 비교하면 \cref{tab:svd-ata}와 같다. 특이값 분해를 직접 계산하면 가장 작은 특이값도 유효숫자 일곱 자리가 맞지만, $A^{\mathsf T}A$의 고유값으로 구하면 $10^{-9}$ 아래의 특이값은 한 자리도 맞지 않는다.

\begin{table}[t]
\begin{sdtable}{colspec={ccc}, row{1}={font=\small\bfseries}}
\toprule
참 특이값 & 특이값 분해의 상대오차 & $\sqrt{\lambda(A^{\mathsf T}A)}$의 상대오차 \\
\midrule
$1.0$ & $1.1\times10^{-16}$ & $2.2\times10^{-16}$ \\
$1.4\times10^{-3}$ & $9.4\times10^{-16}$ & $1.1\times10^{-12}$ \\
$1.9\times10^{-6}$ & $8.1\times10^{-13}$ & $8.0\times10^{-6}$ \\
$2.7\times10^{-9}$ & $2.2\times10^{-9}$ & $1.9$ \\
$1.0\times10^{-10}$ & $4.0\times10^{-8}$ & $53$ \\
\bottomrule
\end{sdtable}
\caption{특이값이 $1$부터 $10^{-10}$까지인 $50\times8$ 행렬에서 두 방법으로 구한 특이값의 상대오차. numpy의 \code{svd}와 \code{eigvalsh}로 binary64에서 계산했다.}
\label{tab:svd-ata}
\end{table}

\section{이중대각화}
\label{sec:bidiagonalization}

23장에서 고유값 알고리즘이 먼저 Hessenberg 형식으로 줄였듯이, 특이값 알고리즘은 먼저 $A$를 \term{이중대각}{bidiagonal} 행렬, 곧 대각과 그 바로 위 대각만 $0$이 아닌 행렬로 줄인다. 왼쪽에서 Householder 반사로 첫 열의 대각 아래를 소거하고, 오른쪽에서 다른 반사로 첫 행의 둘째 원소 오른쪽을 소거하는 일을 번갈아 되풀이하면
\[
A=U_1BV_1^{\mathsf T},\qquad B=\begin{pmatrix}d_1&e_1&&\\&d_2&\ddots&\\&&\ddots&e_{n-1}\\&&&d_n\end{pmatrix}
\]
을 얻는다\cite{golub1965kahan}. 양쪽에 서로 다른 직교 변환을 쓸 수 있는 것은 특이값이 $U^{\mathsf T}AV$에서 변하지 않기 때문이다. 고유값 문제에서는 상사 변환이어야 해서 양쪽이 같아야 했고 그래서 Hessenberg까지만 갈 수 있었지만, 특이값에서는 한 걸음 더 나아간 이중대각까지 유한 번의 단계로 갈 수 있다. 연산은 $4mn^2-4n^3/3$번이다. 행이 열보다 훨씬 많으면 먼저 21장의 QR 분해로 $A=QR$을 구하고 $n\times n$인 $R$을 이중대각화하는 것이 더 싸다. 이중대각화도 21장의 compact WY 표현으로 블록 알고리즘이 되지만, 반사의 절반이 행렬--벡터 곱으로 남아서 QR만큼 빠르지는 않다.

\section{이중대각 행렬의 특이값}
\label{sec:bidiagonal-svd}

\subsection{암시적 QR}

이중대각 행렬 $B$의 특이값은 대칭 삼중대각 행렬 $T=B^{\mathsf T}B$의 고유값의 제곱근이다. $T$에 이동 QR 알고리즘을 적용하는 일을, $T$를 만들지 않고 $B$ 위에서 할 수 있다. 23장의 암시적 Q 정리에 따르면 한 걸음은 첫 열만으로 정해지므로, $T-\mu I$의 첫 열을 소거하는 Givens 회전 하나를 $B$의 오른쪽에 적용해 불룩한 원소를 만들고, 왼쪽과 오른쪽에서 Givens 회전을 번갈아 적용해 그 원소를 대각을 따라 아래로 밀어내면 된다\cite{golub1970reinsch}. 걸음 하나가 $O(n)$이고, 23장에서 본 Wilkinson 이동으로 수렴이 빠르다. 특이벡터가 필요하면 회전들을 $U$와 $V$에 누적한다.

\subsection{상대 정확도}

이중대각 행렬에는 일반 행렬에 없는 성질이 있다. 특이값이 원소들의 작은 상대 변화에 대해 \textbf{상대적으로} 잘 정해진다.

\begin{theorem}[Demmel--Kahan]
\label{thm:bidiag-relative}
$n\times n$ 이중대각 행렬 $B$의 각 원소를 상대적으로 $\epsilon$ 이하만큼 바꾼 행렬을 $\tilde B$라 하면, 대응하는 특이값은 $|\tilde\sigma_i-\sigma_i|\le\bigl((1+\epsilon)^{2n-1}-1\bigr)\sigma_i\approx(2n-1)\epsilon\,\sigma_i$를 만족한다.
\end{theorem}

\begin{proof}
원소 하나를 $1+\epsilon_j$배 하는 것은, 이중대각 꼴 때문에 대각 행렬을 왼쪽과 오른쪽에 곱하는 것으로 나타낼 수 있다. 예컨대 $d_k$만 바꾸려면 앞쪽 행들과 열들에 적당한 배수를 나누어 곱하면 된다. 그래서 $\tilde B=D_1BD_2$이고, $D_1$과 $D_2$는 대각 원소가 $1$에서 상대적으로 $(1+\epsilon)^{2n-1}-1$ 이내인 대각 행렬이다. 그런 $D$에 대해 $\sigma_i(D_1BD_2)$는 $\sigma_i(B)$와 $\|D_1\|\|D_2\|$배 및 $\|D_1^{-1}\|^{-1}\|D_2^{-1}\|^{-1}$배 사이에 있다. 이것은 특이값의 최소최대 특징짓기에서 바로 나온다.
\end{proof}

곧 $B$의 원소를 반올림 수준으로 바꾸어도 가장 작은 특이값까지 모든 특이값이 상대적으로 $nu$ 정도만 움직인다. 일반 행렬에서는 작은 특이값의 상대오차가 $u\sigma_1/\sigma_i$로 커지는 것과 대조적이다. 이 성질을 살리려면 알고리즘이 원소를 상대적으로 작게만 바꾸어야 하는데, 이동 QR은 이동을 빼는 순간 이 성질을 잃는다. Demmel과 Kahan은 이동이 $0$인 QR 걸음을 원소의 상대 정확도를 지키도록 다시 짜서, 이중대각 행렬의 모든 특이값을 높은 상대 정확도로 구하는 방법을 만들었다\cite{demmel1990kahan}. \cref{tab:svd-ata}에서 특이값 분해의 오차가 가장 작은 특이값에서도 $10^{-8}$에 그친 것은, 이중대각화 단계가 $u\sigma_1$ 정도의 절대오차를 만들 뿐이고 그 뒤의 이중대각 단계가 정확도를 거의 잃지 않기 때문이다.

\subsection{dqds}

특이값만 필요하면 더 빠르고 더 정확한 방법이 있다. Rutishauser의 몫--차 알고리즘\cite{rutishauser1954qd}을 미분 꼴로 다시 쓴 \term{dqds}{differential quotient-difference with shifts} 알고리즘은 이중대각 행렬의 원소의 제곱을 직접 갱신하며, 뺄셈 없이 곱셈과 나눗셈과 덧셈만으로 한 걸음을 하므로 상대 정확도를 잃지 않는다\cite{fernando1994dqds}. 이동을 써서 빠르게 수렴하고, 특이벡터를 누적하지 않으므로 걸음 하나가 아주 싸다. LAPACK은 특이값만 구할 때 이 알고리즘을 쓴다.

\section{분할 정복}
\label{sec:svd-dc}

특이벡터까지 모두 구할 때는 23장의 대칭 삼중대각 분할 정복과 같은 생각을 이중대각 행렬에 쓴다\cite{gu1995svd}. 이중대각 행렬을 한가운데에서 두 블록과 계수 하나짜리 보정으로 나누고, 두 블록의 특이값 분해를 재귀적으로 구한 뒤, 계수 하나만큼 바뀐 문제의 특이값을 영년 방정식으로 푼다. 마지막에 특이벡터를 합치는 곱이 3층 BLAS의 행렬 곱이므로, 큰 행렬에서는 암시적 QR보다 몇 배 빠르다. LAPACK의 \code{gesdd}가 이것이고, \code{gesvd}는 암시적 QR이다.

\begin{instardust}[\code{docs/design/linalg.md} §10.4, §10.5]
stardust의 \code{svd}는 \code{gesdd}를, 특이값만 구하는 \code{svdvals}와 그 위의 \code{matrix\_rank}, \code{cond}, 2-노름은 \code{gesvd}의 특이값 전용 모드를 쓴다. 특이값 전용 모드는 $1.6$--$2$배 빠르지만 그 값이 \code{svd}가 주는 특이값과 $10^{-15}$ 수준에서 다를 수 있으므로, 특이값만 쓰는 함수들이 모두 같은 도우미를 거쳐 서로 비트까지 같은 값을 쓰게 했다. \code{svd}의 특이값은 numpy와 같지만 특이벡터는 다를 수 있는데, 설계 문서에 따르면 \code{gesdd}가 작업 공간의 크기를 보고 내부 경로를 고르고 numpy와 stardust가 그 크기를 다르게 정하기 때문이다. 유사역행렬 \code{pinv}에는 10장의 거듭제곱 조정이 등장한다. 남길 특이값이 $2^{-1022}$ 아래이면 그 역수가 넘쳐 무한대가 되고, 특이벡터의 정확한 $0$과 곱해져 NaN이 된다. 그래서 그런 항만 특이값에 $2^{600}$을 곱해 따로 더한 뒤 마지막에 $2^{600}$을 곱해 되돌린다. $2$의 거듭제곱이라 이 조정은 정확하다.
\end{instardust}

\begin{aside}[특이값 분해의 오랜 역사]
특이값 분해는 1870년대에 Beltrami와 Jordan이 이차 형식을 연구하며 서로 독립적으로 발견했고, 20세기 초 Schmidt는 적분 연산자에 대해, Eckart와 Young은 1936년 최선 저랭크 근사 정리(4장)로 통계학과 심리측정학에 쓰일 길을 열었다. 그러나 그것을 안정하게 계산하는 방법은 1965년 Golub과 Kahan의 이중대각화가 나오고서야 생겼고, 1970년의 Golub--Reinsch 알고리즘이 오늘날 소프트웨어의 원형이 되었다\cite{stewart1993svdhistory}. 이론이 나온 지 거의 한 세기 만에 계산이 따라잡은 것이다.
\end{aside}

\section{맺음}

특이값 분해의 계산은 이 부의 앞선 장들이 모두 합쳐지는 곳이다. 21장의 Householder 반사로 이중대각화하고, 23장의 암시적 이동 QR이나 분할 정복으로 이중대각 행렬의 특이값을 구하며, 22장에서 본 정규방정식의 함정을 $A^{\mathsf T}A$를 만들지 않음으로써 피한다. 이중대각 행렬의 특이값이 원소에 대해 상대적으로 잘 정해진다는 성질은, 알고리즘을 그 구조에 맞추면 일반적인 오차 해석이 약속하는 것보다 훨씬 정확한 결과를 얻을 수 있음을 보여 준다. 다음 장에서는 행렬의 함수, 특히 행렬 지수를 계산하는 방법과 Sylvester 방정식을 다룬다.

\chapterreferences
